At $n=4$ we have $N=81$, $s\ge4$ means $81\mid\ell^2-1$, i.e.\ $\ell\equiv\pm1\pmod{81}$, so $r=4$, $c=54$, and $f=1$ for $\ell\equiv1$, $f=2$ for $\ell\equiv-1\pmod{81}$; $L^{\mathrm{rel}}_{3,4}=9\log\frac{3^{40/81}+\sqrt{3^{80/81}+4}}2=7.0107178918636975\ldots$. The certificate \artPcert{} is the outward hull, widened to twice its half-width, of two independent $4000$-bit ball enclosures of every quantity below: the producer's (the $81$ logarithms $\log|\sigma^t\eta_4|$ are balls, the $54\times54$ Gram matrix and its determinant are ball operations, and independently the discrete Fourier transform of the log vector and the character product of Theorem~\ref{thm:rank3} Step~4; the two enclosures of $D_4^{(4)}$ have radius below $10^{-1000}$ and their quotient encloses $1$) and the read-only replay's \artPreplay{}, which rebuilds the same quantities from the definitions by Gram--Schmidt orthogonalisation and a table of powers of $e^{-2\pi i/N}$; the certificate is generated by \artPgen{}, and the replay requires every certified interval to contain its own recomputation (coverage checker \artPcover{}). It encloses
\begin{align*} C^{(3)}_{4,4}&=\frac{(2/\pi)^{27}\,\Gamma(29)\,D_4^{(4)}}{(L^{\mathrm{rel}}_{3,4})^{54}}\in[1.072749779\times10^{33},1.072749780\times10^{33}],\\ \sqrt{C^{(3)}_{4,4}}&\in[3.275285910\times10^{16},3.275285911\times10^{16}]. \end{align*}
Hence $\ell\equiv1\pmod{81}$ and $\ell>1.0728\times10^{33}$ give $\ell^{f}=\ell>C^{(3)}_{4,4}$, and $\ell\equiv-1\pmod{81}$ and $\ell>3.2753\times10^{16}$ give $\ell^{f}=\ell^2>C^{(3)}_{4,4}$; Theorem~\ref{thm:P3} applies. The same certificate gives the rows $n\le5$ of the comparison table in the $\Z_3$ section of the paper. For the comparison: the uniform-in-$n$ bound of Morisawa and Okazaki, \cite{MO16} Thm.~A, which for $p=3$, $s=4$ reads $\ell>G(3,4,f)$ with $G(3,s,f)=\big((\sqrt{2\pi}/(3^{3/4}\log\frac{3^{40/81}+\sqrt{3^{80/81}+4}}2))^{c}\,((c+2)/2)!\big)^{1/f}$, $c=2\cdot3^{s-1}$, i.e.\ $G(3,4,1)\ge3.7111\times10^{37}$ and $G(3,4,2)\ge6.0919\times10^{18}$ (the same certificate encloses $G$; its value $G(3,3,2)=42407.49\ldots$ agrees with the table of \cite{MO16} Example~1.6, ``$4.3\times10^4$''). In the class $\ell\equiv1\pmod{81}$ with $81\,\|\,\ell-1$ (the set $D(3,4,1)$ of \cite{MO16}) the present threshold is smaller by the factor $G(3,4,1)/C^{(3)}_{4,4}\ge3.4594\times10^{4}$, and in $D(3,4,2)$ ($\ell\equiv-1\pmod{81}$, $81\,\|\,\ell+1$) by $G(3,4,2)/\sqrt{C^{(3)}_{4,4}}\ge1.8599\times10^{2}$. The comparison convention is that of the paper's comparison section (layer-fixed versus uniform in $n$): \cite{MO16} is uniform in $n$ (and its constant for $s\ge5$ is $G(3,s,f)$, larger), the present bound is for the layer $n=4$ only; \cite{MO13} Thm.~0.3 ($2^{c/2}c!=3.098\times10^{79}$ at $s=4$) and Thm.~A ($G_1(3,4,f)^f=7.63\times10^{101}$) are listed in that table for the record, not as the comparator.
